\begin{textsample}{POS Dim 3 – persona\_gemini – Score 18.00 – t850\_gemini.txt}  \label{ex:f3_pos_044}
Since 2007, 7.1 \textbf{billion} smartphones have been produced, \textbf{using} a complex mix of over 60 different \textbf{elements}.
The mining and processing of these \textbf{materials} have significant environmental and social impacts.
The \textbf{energy} consumed in the manufacturing of these devices has reached 968 TWh, an \textbf{amount} \textbf{equivalent} to the annual \textbf{energy} \textbf{supply} of India.
This \textbf{production} has led to a massive e-waste \textbf{problem}.
In 2014, smartphones contributed to 3 million metric \textbf{tons} of e-waste, yet only 16 % of this was recycled.
The \textbf{design} of the devices contributes to this issue, as most \textbf{models} lack easily replaceable batteries.
This \textbf{design} choice results in short device lifespans, which are about two years in the US. \textbf{Recycling} is often inefficient, and the recall of the Samsung Galaxy Note 7 highlights the \textbf{scale} of the \textbf{waste}.
A \textbf{shift} to a circular \textbf{model} is necessary.
We must call on \textbf{companies} like Samsung to improve the repairability of their devices and commit to \textbf{better} \textbf{recycling}.
Petitions are being organized to urge them to take these steps.

% matched lemmas: amount, billion, company, design, element, energy, equivalent, good, material, model, problem, production, recycling, scale, shift, supply, ton, use, waste
\end{textsample}
